\documentclass[12pt]{article}
\usepackage[utf8]{inputenc}
\usepackage[T1]{fontenc}
\usepackage{hyperref}
\title{Woleaian Language Resources for K-12 Curriculum Development}
\author{Compiled by Claude Sonnet 5.5 (claude-sonnet-5-5) for Andr\'as Kornai}
\date{}
\begin{document}
\maketitle

\section{Introduction}
Woleaian (ISO 639-3: woe; Glottocode: wole1240) is an Austronesian language of Micronesia with about 1,600 speakers. It is not known to be taught in schools. The searches turned up reference works (a dictionary and a grammar) but no primers, textbooks or school programmes.

\section{Primary Educational Resources}
No Woleaian primers, textbooks or classroom materials turned up in the searches. Any curriculum will have to be built from the reference works below.

\section{Linguistic and Reference Materials}
\begin{itemize}
\item \emph{Woleaian-English Dictionary} (Ho-min Sohn and Anthony Tawerilmang, 1976, University of Hawai'i Press). It has about 6,200 Woleaian entries and an English-Woleaian finder list of about 4,000 entries. The UH Press Manifold edition is under a CC BY-NC-SA 4.0 licence. \url{https://manifold.uhpress.hawaii.edu/projects/woleaian-english-dictionary}
\item ScholarSpace record for the dictionary: \url{https://scholarspace.manoa.hawaii.edu/items/51651da3-5029-4b86-a343-aca8c94167c4/full}
\item Habele Institute reference entry for the dictionary: \url{https://habeleinstitute.org/wiki/Zotero_reference:X6GLC5C4}
\item \emph{Woleaian Reference Grammar} (Ho-min Sohn, 1975, PALI Language Texts: Micronesia). The record is available through DOAB: \url{https://oalib-perpustakaan.upi.edu/Record/doab_20_500_12854_114450/Holdings}
\item Sohn's 1984 \emph{An Orthographic Design for Woleaian} and Tawerilmang's 1984 \emph{Proto-Oceanic Reflexes in Woleaian} are listed on tekinged.com. No free copy was seen, so they are mentioned here only.
\end{itemize}

\section{Religious and Community Resources}
No Woleaian scripture or community-organisation resources were found in the searches.

\section{Digital and Audio Resources}
\begin{itemize}
\item tekinged.com hosts the dictionary and grammar, plus a partial online Woleaian-English dictionary. The site says it is not actively maintained and may be incomplete. \url{https://tekinged.com/books/woleaian_dict.php}
\end{itemize}
No audio or video resources were found.

\section{Curriculum Development Context}
The language is stable and used by the whole community, but the searches found no school materials. The 1976 dictionary omits archaic words that younger speakers no longer use. Curriculum work would start from the openly licensed dictionary and the grammar, with new primers and readers written with community speakers.

\end{document}